\scanpage{67}
% ll. 1–3
on normalise \inserted{$\lambda_0,\lambda_1$} par le choix de racines primitives $3^{\mathrm e},4^{\mathrm e}$ et $5^{\mathrm e}$ de $1$ (par exemple $\exp 2i\pi/\nu$ quand on est sur $\mathbb C$), et ceci pour deux raisons :

% ll. 4–10
1) Cas où $\lambda_0$ ou $\lambda_1$ \unsure{estt|est} égal à $-1$, i.e. $\nu_0$ ou $\nu_1$ égal 2 --- cela introduit (pour $K$ déjà fixé) une ambiguïté d'ordre 2 :

\diagram{$\begin{array}{l@{\quad}l@{\quad}l}
\text{Cas}&(3,2,3)\quad(2,3,3)&(\inserted{\text{cas}}\ \text{tétraédraux})\\
&(3,2,4)(2,4,3)(2,3,4)(4,2,3)&(\text{octaédraux})\\
&\left\{\begin{array}{l}(3,2,\underline{5})(2,5,3)(2,3,\underline{5})(5,2,3)\\
(5,2,5')(2,5,5')\end{array}\right.&\text{(icosaédraux}
\end{array}$}

% ll. 11–14
2) Cas où $\nu=5$, et où on peut avoir le choix entre $K=\alpha$ et $K=\alpha'$ (cas icosaédral)
\[
\begin{aligned}
&(3,2,5)(2,3,5)\qquad\text{(qui figurent déjà ds liste 1)}\\
&[(3,2,5')\ (2,3,5')]
\end{aligned}
\]
\deleted{\unsure{Si}}

% ll. 15–22
Finalement, compte tenu du choix des racines primitives $3^\circ\ 4^\circ$ et $5^\circ$ de $1$ pour normaliser $\lambda_0,\lambda_1$, on trouve les possibilités suivantes --- où on souligne les possibilités numériques qui donnent lieu à \uline{deux} hyperpolygônes, en interprétant comme d'habitude $5,5'$ comme signifiant $K=\alpha$ resp. $K=\alpha'$ :
% archive sheet number: 66
